% Folder 83, batch 11 (archivists' pages 201-220).
% Pass: Opus 5.5 (claude-opus-5-5), 2026-10-10 — first pass, unchecked against the pages by a human.
%
% Pages 201, 203, 207, 209, 211, 213, 215, 217 and 219 are typescript pages of
% Récoltes et Semailles (with his blue-ink corrections), unrelated to the
% mathematics in hand; page 205 is blank. They are not transcribed.
\documentclass[11pt,a4paper]{article}
\input{../preamble/grothendieck}

\folder{83}
\batch{11}
\pages{201}{220}
\foldertitle{[Icosaèdres, bi-icosaèdres et algèbres d'Azumaya de rang 4] : notes manuscrites (1982-1983, 1986, s.d.).}
\dating{1982-1986}
\watermark{Édition de démonstration}

\begin{document}

\section{Relations sur les arêtes d'un cube}

\page{202}
\note{Les calculs de ce lot portent sur un cube dont les sommets, arêtes et faces sont indexés ; le contexte (traces de matrices $X_i$, p. 204) est vraisemblablement celui des pages précédentes du dossier, hors de ce lot.}

\drawing{Un cube en perspective. Sommets nommés $A$, $B'$, $C'$, $A'$, $B$, $C$ ; deux sommets opposés marqués d'un point plein. Arêtes nommées $a$ (verticale issue de $A$), $b$, $c$, $a'$, $b'$, $c'$ ; les arêtes $AB'$, $C'$–(sommet marqué) et $a$ sont repassées en trait gras ; un « $W$ » est écrit près de $B$.}

\marginal{Trop lourd}\note{encerclé, en haut à droite de la page.}

\[
\Sigma_a\quad \Sigma_b\quad \Sigma_c\quad \Sigma_{AB'}\quad \Sigma_{BC'}\quad \Sigma_{CB'}\quad \Sigma_{CA'}\quad \Sigma_{AC'}
\]
\[
\Sigma_{a'}\quad \Sigma_{b'}\quad \Sigma_{c'}
\]
\note{Dans $\Sigma_{CB'}$ de la première ligne, l'indice $C$ est écrit en surcharge d'une autre lettre (vraisemblablement $A$) ; même surcharge dans les relations ci-dessous.}

\[
\left\{
\begin{aligned}
\Sigma_a + \Sigma_{c'} &= \Sigma_{AB'} + \Sigma_{BC'}\\
\Sigma_a + \Sigma_{b'} &= \Sigma_{AC'} + \Sigma_{CB'}\\
\Sigma_b + \Sigma_{a'} &= \Sigma_{BC'} + \Sigma_{CA'}\\
\Sigma_b + \Sigma_{c'} &= \Sigma_{BA'} + \Sigma_{AC'}\\
\Sigma_c + \Sigma_{b'} &= \Sigma_{CA'} + \Sigma_{AB'}\\
\Sigma_c + \Sigma_{a'} &= \Sigma_{CB'} + \Sigma_{BA'}
\end{aligned}
\right.
\]
\note{Des accolades relient deux à deux, à droite, des lignes de ce premier groupe.}

\[
\left[
\begin{aligned}
\Sigma_a + \Sigma_{CA'} &= \Sigma_b + \Sigma_{CB'}\,.\\
\Sigma_{a'} + \Sigma_{AC'} &= \Sigma_{b'} + \Sigma_{BC'}\\
\Sigma_b + \Sigma_{AB'} &= \Sigma_c + \Sigma_{AC'}\,.\\
\Sigma_{b'} + \Sigma_{BA'} &= \Sigma_{c'} + \Sigma_{CA'}\,.\\
\Sigma_c + \Sigma_{BC'} &= \Sigma_a + \Sigma_{BA'}\,.\\
\Sigma_{c'} + \Sigma_{CB'} &= \Sigma_{a'} + \Sigma_{AB'}
\end{aligned}
\right.
\]
\note{Les indices $AC'$ (deuxième ligne), $BA'$ (quatrième) et $CB'$ (dernière) sont écrits en surcharge ; le $c'$ de la quatrième ligne aussi. La relation de la deuxième ligne est recopiée à droite : $\Sigma_{a'} + \Sigma_{AC'} = \Sigma_{b'} + \Sigma_{BC'}$.}

En marge, à droite du premier groupe :

$\exists\, \sigma$ avec $\sigma_{\underline{a}} + \sigma_{\underline{b}} + \sigma_{\underline{c}} = \sigma$ $\forall$ triple d'arêtes formant un « demi hexagone »

\drawing{Petit cube au crayon, sur lequel trois arêtes deux à deux non coplanaires $\underline{a}$, $\underline{b}$, $\underline{c}$ sont repassées à l'encre, formant une ligne brisée.}

\uncertain{Chacune des six dernières relations est conséquence des cinq autres, comme on voit en sommant, ce qui donne} $0=0$ !
\note{lecture d'ensemble d'une phrase rapide, en marge à droite du second groupe ; plusieurs mots sont à peine formés.}

\[
3\bigl(\Sigma_a - \Sigma_b + \Sigma_c + \Sigma_{a'} - \Sigma_{b'} + \Sigma_{c'}\bigr) + \bigl(\Sigma_{CA'} + \Sigma_{AC'} - \cdots\bigr)
\]
\[
= \bigl(\Sigma_a + \Sigma_b - \Sigma_c - \Sigma_{a'} - \Sigma_{b'} - \Sigma_{c'}\bigr) + 3\bigl(\Sigma_{CA'} + \Sigma_{AC'}\bigr)
\]
\note{Les signes à l'intérieur des deux premières parenthèses sont mal formés et leur lecture est incertaine ; la première ligne s'interrompt sur deux tirets.}

\section{Variables et relations : traces sur les sommets, arêtes et faces}

\page{204}
\note{Dans ce lot, une lettre $T$ à hampe doublée, épaissie, désigne des traces de produits de trois ou quatre matrices ; elle est rendue $\mathbb{T}$. Le $T$ simple (p. 206) est conservé tel quel.}

\[
(4) + (4) + (6)\qquad
\boxed{\begin{array}{ll} b_s\ \ c_s & (s\in S)\\ t_a & (a\in A)\end{array}}
\qquad\text{avec}\quad b_s = b_{\bar s},\ c_s = c_{\bar s},\ t_a = t_{\bar a}
\]
\note{Les nombres 4, 4, 6 sont cerclés. $b_s$ et $c_s$ sont écrits en surcharge de $a_s$, $b_s$ ; chacune des trois variables est cerclée.}

\drawing{Un cube en perspective ; trois sommets marqués $s_1$, $s_2$, $s_3$ autour d'un sommet $s$ cerclé ; chaque face visible porte une petite flèche circulaire (orientation).}

Si $\delta(s)$ \add{$\in I$} est la diagonale qui passe par $s$, on a
\[
b_s = \operatorname{Tr} X_s,\qquad c_s = \det X_s .
\]
Si $\delta(a) \in \mathfrak{P}_2(I)$ est l'ensemble des deux diagonales qui passent par les \struck{\ill{}} \uncertain{extrémités} de $a$, on a
\[
t_a = \operatorname{Tr} X_i X_j \;=\; \operatorname{Tr} X_j X_i \quad\text{si } \delta(a) = \{i,j\}.
\]

\[
4 + 4 + (8)\qquad
\boxed{\sigma_s,\ \pi_s,\ \mathbb{T}_s}\qquad (s\in S)
\]

$s$ correspond à un triple $(i,j,k)$ à permutation circulaire près
\[
\mathbb{T}_s = \operatorname{Tr} X_i X_j X_k = \operatorname{Tr} X_j X_k X_i = \operatorname{Tr} X_k X_i X_j
\]
\[
\mathbb{T}_{\bar s} = \operatorname{Tr} X_k X_j X_i = \operatorname{Tr} X_j X_i X_k = \operatorname{Tr} X_i X_k X_j
\]
\note{Le deuxième terme de la première ligne est écrit en surcharge d'un premier essai.}

(4 + 4 relations)\note{encerclé.}
\[
(R_s)\ \sim\ (R_{\bar s})\qquad
\boxed{\begin{aligned}
\sigma_s &= \sigma_{\bar s} = \mathbb{T}_s + \mathbb{T}_{\bar s}\\
\pi_s &= \pi_{\bar s} = \mathbb{T}_s \cdot \mathbb{T}_{\bar s}
\end{aligned}}
\]
\begin{align*}
\sigma_s = \sigma_{\bar s} &= b_i t_{jk} + b_j t_{ki} + b_k t_{ij} - b_i b_j b_k\\
&= b_{s_1} t_{ss_1} + b_{s_2} t_{ss_2} + b_{s_3} t_{ss_3} - b_{s_1} b_{s_2} b_{s_3}
\end{align*}
où $s_1, s_2, s_3$ sont les sommets adjacents à $s$
\begin{align*}
\pi_s = \pi_{\bar s} &= t_{ss_1} t_{ss_2} t_{ss_3}
+ \text{\struck{$(c_{s_1} t_{ss_1} + c_{s_2} t_{ss_2} + c_{s_3} t_{ss_3})$}}\\
&\quad + \bigl[c_{s_1} t_{ss_1}(t_{s_2 s_3} - b_{s_2} b_{s_3}) + c_{s_2} t_{ss_2}(t_{s_3 s_1} - b_{s_3} b_{s_1})\\
&\qquad + c_{s_3} t_{ss_3}(t_{s_1 s_2} - b_{s_1} b_{s_2})\bigr]\\
&\quad + \bigl(b_{ss_1}^2\, c_{ss_2} c_{ss_3} + b_{ss_2}^2\, c_{ss_3} c_{ss_1} + b_{ss_3}^2\, c_{ss_1} c_{ss_2}\bigr)
\end{align*}
\note{Les indices des $t$ entre parenthèses dans le crochet sont surchargés et leur lecture ($t_{s_2s_3}$, etc.) est incertaine ; les indices doubles de la dernière parenthèse sont transcrits tels qu'ils semblent écrits.}

\[
(6)\qquad
\boxed{\begin{array}{ll} \mathbb{T}_f & (f\in F)\\ \sigma_a,\ \pi_a & (a\in A)\end{array}}
\]
Si $f$ correspond à la permutation circulaire $(i,j,k,l)$ de $I$, on a
\[
\text{\struck{$\mathbb{T}_f =$}}\qquad
\mathbb{T}_f = \operatorname{Tr} X_i X_j X_k X_l = \operatorname{Tr} X_j X_k X_l X_i = \operatorname{Tr} X_k X_l X_i X_j = \operatorname{Tr} X_l X_i X_j X_k
\]

Si $f$, $f'$ sont adjacentes le long d'une arête $a = (s,t)$

\drawing{Petit cube : la face $f$ de sommets $s, t, s', t'$ et la face $f'$ de sommets $t, s, t'', s''$, adjacentes le long de l'arête $a$ entre $s$ et $t$ ; arêtes $a'$, $a''$ marquées, ainsi que la diagonale $i$.}

\[
f = (s, t, s', t'),\quad f' = (t, s, t'', s'')
\]
\note{Sous les sommets de $f$ sont inscrits $i, j, k, l$, sous ceux de $f'$ : $j, i, k, l$.}

\struck{\ill{}} \ill{} $\mathbb{T}_f = \operatorname{Tr} X_i X_j (X_k X_l)$, $\mathbb{T}_{f'} = \operatorname{Tr} X_j X_i (X_k X_l)$, on a

\page{206}
$(R_a)$\note{cerclé.} 12 + 12 relations
\[
\boxed{\begin{aligned}
\mathbb{T}_f + \mathbb{T}_{f'} &= \sigma_a\\
\mathbb{T}_f\, \mathbb{T}_{f'} &= \pi_a
\end{aligned}}
\]
avec
\begin{align*}
\sigma_a &= \operatorname{Tr} X_k X_l \operatorname{Tr} X_i X_j + \operatorname{Tr} X_i \operatorname{Tr} X_j X_k X_l + \operatorname{Tr} X_j \operatorname{Tr} X_k X_l X_i\\
&\qquad - \operatorname{Tr} X_i \operatorname{Tr} X_j \operatorname{Tr} X_k X_l\\
&\text{\struck{$= t_{s\ldots}\, t$}}\\
&= t_{a'a''}\, t_a + b_s \mathbb{T}_s + b_t T_t - b_s b_t t_{a',a''}
\end{align*}
\note{Le premier facteur $t_{a'a''}$ est écrit en surcharge et glosé dessous par « $= t_{a''}$ » ; lecture incertaine.}

$\pi_a = \cdots$ expression \add{un peu} \struck{plus} compliquée en \struck{$\operatorname{Tr} X_i$, …} $b_i, c_i, b_j, c_j, t_{kl}, c_k, c_l$,
$t_{ij}$, $\underbrace{T_{jkl}}_{T_s}$, $\underbrace{T_{kli}}_{T_t}$

\[
\overbrace{4\quad 4\quad 6}^{14}\qquad \overbrace{8\quad 6}^{14}
\]
\[
28 \text{ variables}\quad b_s,\ c_s,\ t_a,\ \mathbb{T}_s,\ \mathbb{T}_f
\]

32 relations !

\section{Paramétrage des arêtes}

\page{208}
\drawing{Deux cubes en perspective. Sur le petit, en haut à gauche, les arêtes portent $a$, $b$, $c$, $d$, $a'$, $b'$ et, en dessous, $a+\alpha$, $b+\alpha$, $c+\alpha$, $d+\alpha$. Sur le grand, les arêtes portent, cerclées, $a$, $b$, $c$, $d$, $u$, et $a+\alpha = a'$ (cerclé $a'$), puis $b+\alpha$, $c+\alpha$, $d+\alpha$, $u+c-a$ ; sur une face : $v = u - (d-b)$ (signe surchargé).}

\[
v + (c-a) = u - (d-b) + (c-a)
\]
\[
\boxed{d - b = u - v} = d - b = u - v
\]
i.e.
\[
d + v = b + u,\qquad d - u = b - v
\]
\note{La seconde égalité ne découle pas de la première ; transcrite telle quelle.}

$a, b, c, d, \alpha, u$

ou $a + \alpha = a'$ !

\page{210}
$a, b, c, d, a'$ \struck{\ill{}} ; $u$
\[
\alpha = a' - a
\]
\uncertain{Sommes} \ill{} sur les arêtes des six faces\note{en haut, au centre de la page ; lecture incertaine.}

\begin{align*}
S_1 &= a + b + c + d\\
S'_1 &= a + b + c + d + 4(a'-a) = -3a + b + c + d + 4a'\\
S_2 &= a + u + a' + u + (d-b) = a - b + \cdot + d + a' + 2u\\
S'_2 &= c + u + c - a + c + (a'-a) + u - d + b + c - a\\
&= -3a + b + 4c - d + a' + 2u\\
S_3 &= u + d + a' - a + u + c - a + d = -2a + \cdot + c + 2d + a' + 2u\\
S'_3 &= b + u - d + b + b + a' - a + u - d + b + c - a\\
&= -2a + 4b + c - 2d + a' + 2u
\end{align*}
\note{Dans $S'_1$, le terme « $+\,d$ » est ajouté au-dessus de la ligne. Dans $S'_2$, « $(a'-a)$ » est écrit au-dessus d'un terme surchargé ; dans la ligne finale de $S'_2$, le coefficient $4$ de $c$ est en surcharge. Les termes des sommes sont soulignés à mesure qu'ils sont comptés.}

\[
\begin{array}{cccccc}
a & b & c & d & a' & u
\end{array}
\]
\[
\begin{pmatrix}
1 & 1 & 1 & 1 & 0 & 0\\
-3 & 1 & 1 & 1 & 4 & 0\\
1 & -1 & 0 & 1 & 1 & 2\\
-3 & 1 & 4 & -1 & 1 & 2\\
-2 & 0 & 1 & 2 & 1 & 2\\
-2 & 4 & 1 & -2 & 1 & 2
\end{pmatrix}
\]

\[
\Delta = \det\begin{pmatrix}
1 & 1 & 1 & 4 & 0\\
-1 & 0 & 1 & 1 & 2\\
1 & 4 & -1 & 1 & 2\\
0 & 1 & 2 & 1 & 2\\
4 & 1 & -2 & 1 & 2
\end{pmatrix}
- \det\begin{pmatrix}
-3 & 1 & 1 & 4 & 0\\
1 & 0 & 1 & 1 & 2\\
-3 & 4 & -1 & 1 & 2\\
-2 & 1 & 2 & 1 & 2\\
-2 & 1 & -2 & 1 & 2
\end{pmatrix}
\]
\[
+ \det\begin{pmatrix}
-3 & 1 & 1 & 4 & 0\\
1 & -1 & 1 & 1 & 2\\
-3 & 1 & -1 & 1 & 2\\
-2 & 0 & 2 & 1 & 2\\
-2 & 4 & -2 & 1 & 2
\end{pmatrix}
- \det\begin{pmatrix}
-3 & 1 & 1 & 4 & 0\\
1 & -1 & 0 & 1 & 2\\
-3 & 1 & 4 & 1 & 2\\
-2 & 0 & 1 & 1 & 2\\
-2 & 4 & 1 & 1 & 2
\end{pmatrix}
\]
\note{Développement selon la première ligne ; le calcul s'arrête là.}

\section{Les paramètres $\delta$, $\delta'$, $\delta''$}

\page{212}
\drawing{Un cube en perspective, sommets $A$, $B'$, $C$, $B$, $A'$ nommés, un sommet marqué d'un point plein à l'origine des arêtes $a$, $b$, $c$, un autre sommet marqué d'un point. Arêtes nommées $a$, $b$, $c$, $a'$, $b'$, $c'$, $\alpha$, $\alpha'$, $\beta$, $\beta'$, $\gamma$, $\gamma'$ ; l'arête $\beta$ ($AB'$) est doublée ; un petit arc près du sommet marqué.}

\[
\gamma' - c = a' - \beta' = c' - \beta = \alpha - a \;\Big|\; = \delta'
\]
\[
\left\{
\begin{aligned}
\gamma' &= \delta' + c\\
\beta' &= -\delta' + a'\\
\beta &= -\delta' + c'\\
\alpha &= \delta' + a
\end{aligned}
\right.
\qquad \text{reste pour compte } \underline{\gamma, \alpha'}
\]

\[
b - \gamma = \underline{\alpha - a'} = \alpha' - b' = a - \beta' \;\Big|\; = -\delta''
\]
\[
\left\{
\begin{aligned}
\gamma &= \delta'' + b\\
\alpha' &= -\delta'' + b'
\end{aligned}
\right.
\qquad
\begin{aligned}
\alpha &= -\delta'' + a'\\
\beta' &= \delta'' + a
\end{aligned}
\]
\note{Dans la première ligne de l'accolade, un premier signe est surchargé ; de même devant $\delta''$ dans la dernière ligne.}

\[
\boxed{-\delta'' + a' = \delta' + a}\qquad \delta' + \delta'' = a' - a
\]
\note{La dernière relation est cerclée.}
\[
\boxed{+\delta'' + a = -\delta' + a'}
\]

\page{214}
\[
b - \alpha' = \gamma' - c' = \gamma - b' = c - \beta = -\delta
\]
\begin{align*}
\alpha' &= \delta + b\\
\gamma' &= -\delta + c'\\
\gamma &= -\delta + b'\\
\beta &= \delta + c
\end{align*}

\[
\boxed{\text{\struck{$\delta + b = -\delta'' + b'$}}}\qquad
\boxed{-\delta + c' = \delta' + c}
\]
\[
\boxed{\text{\struck{$-\delta + b' = \delta'' + b$}}}\qquad
\boxed{\delta + c = -\delta' + c'}
\]
\note{Les deux encadrés de gauche sont barrés d'un trait oblique ; dans le deuxième, $c'$ et $\delta'$ sont en surcharge.}

En marge, cerclé :
\[
\delta + \delta'' = b' - b,\qquad \delta + \delta' = c' - c
\]

\begin{align*}
\delta' + \delta'' &= a' - a\\
\delta'' + \delta &= b' - b\\
\delta + \delta' &= c' - c
\end{align*}
\[
\delta + \delta' + \delta'' = \tfrac12\bigl(a' + b' + c' - (a + b + c)\bigr)
= \tfrac12\bigl(\underbrace{(a'-a)}_{\varepsilon_1} + \underbrace{(b'-b)}_{\varepsilon_2} + \underbrace{(c'-c)}_{\varepsilon_2}\bigr)
\]
\note{Le troisième terme est indexé $\varepsilon_2$ sur la page ; l'encadré qui suit l'utilise comme $\varepsilon_3$.}
\[
\boxed{\left\{
\begin{aligned}
\delta &= \tfrac12(-\varepsilon_1 + \varepsilon_2 + \varepsilon_3)\\
\delta' &= \tfrac12(\varepsilon_1 - \varepsilon_2 + \varepsilon_3)\\
\delta'' &= \tfrac12(\varepsilon_1 + \varepsilon_2 - \varepsilon_3)
\end{aligned}
\right.}
\]

\page{216}
\note{Ci-dessous, $\bullet$ marque un signe surchargé d'une tache d'encre, illisible.}
\begin{align*}
\alpha &= \sigma_{BC'} = \tfrac12\bigl((a'+a) \bullet (b'-b) + (c'-c)\bigr)\\
\beta &= \sigma_{AB'} = \tfrac12\bigl(-(a' \bullet a) + (b'-b) + (c'+c)\bigr)\\
\gamma &= \sigma_{CA'} = \tfrac12\bigl((a'-a) + (b'+b) - (c'-c)\bigr)\\
\alpha' &= \sigma_{AC'} = \tfrac12\bigl(-(a'-a) + (b'+b) + (c'-c)\bigr)\\
\beta' &= \sigma_{CB'} = \tfrac12\bigl((a'+a) + (b'-b) - (c'-c)\bigr)\\
\gamma' &= \sigma_{BA'} = \tfrac12\bigl((a'-a) \bullet (b'-b) + (c'+c)\bigr)
\end{align*}
\note{Dans chaque ligne, la parenthèse où figure une somme $x'+x$ est soulignée d'un trait ondulé. Des flèches courbes, à droite, relient $\alpha$, $\beta$ et $\gamma$ entre elles, puis $\alpha'$, $\beta'$ et $\gamma'$ ; plusieurs signes ($a'+a$ dans $\beta'$, etc.) sont surchargés et de lecture incertaine.}

En haut à droite :
\[
-c' + c + 2c',\qquad -b' + b,\qquad -a' + a,\qquad c' - c
\]

\uncertain{Sommation des $\sigma_a$ sur les arêtes d'une face}
\begin{align*}
\varphi_a &= (a'-a) + 2(b+c)\\
\varphi_b &= (b'-b) + 2(a+c)\\
\varphi_c &= (c'-c) + 2(a+b)
\end{align*}
\note{Dans la première ligne, un début d'expression illisible est raturé devant $(a'-a)$.}

$a - a' + (b'+c')$
\begin{align*}
\varphi_{a'} &= (b'+c') \bullet (a'-a) + b' + c' = 2(b'+c') - (a'-a)\\
\varphi_{b'} &= a' + c' + a' - (b'-b) + c' = 2(c'+a') - (b'-b)\\
\varphi_{c'} &= (a'+b') + a' + b' - (c'-c) = 2(a'+b') - (c'-c)
\end{align*}
\note{Sous ces lignes, un passage de deux lignes est entièrement raturé de traits en zigzag, terminé par un « ! » ; des flèches en montent vers les membres de droite.}
\struck{\ill{}} !

\page{218}
\[
\begin{array}{c}
\begin{array}{cccccc} a & b & c & a' & b' & c' \end{array}\\
\det\left(\begin{array}{ccc|ccc}
-1 & 2 & 2 & 1 & 0 & 0\\
2 & -1 & 2 & 0 & 1 & 0\\
2 & 2 & -1 & 0 & 0 & 1\\
\hline
1 & 0 & 0 & -1 & 2 & 2\\
0 & 1 & 0 & 2 & -1 & 2\\
0 & 0 & 1 & 2 & 2 & -1
\end{array}\right)
\end{array}
= \det\begin{pmatrix} 27 & 1\\ 1 & 27\end{pmatrix}
= 27^2 - 1 \in \mathbb{Q}^*
\]
\[
= (3^3)^2 - 1 = 3^6 - 1 = 728 = 8\cdot 91 = \boxed{2^3\cdot 7\cdot 13}
\]
\uncertain{premier à} $3$\note{souligné, sous la matrice $2\times2$.}

\[
-1 + 8 + 8 - (-4 - 4 - 4) = 16 + 12 - 1 = 27
\]
\note{Calcul du déterminant du bloc $3\times3$ ; « $16$ » est écrit en surcharge.}

\struck{\ill{}}
\[
\begin{array}{r}
27\\
27\\
\hline
189\\
54\phantom{0}\\
\hline
729
\end{array}
\]

Ok

\page{220}
\begin{align*}
A &= a + a' = \lambda\\
B &= b + b' = \mu\\
C &= c + c' = \nu
\end{align*}
\struck{$\beta$}
\begin{align*}
A' &= a' + a\\
B' &= b' + b\\
C' &= c' + c
\end{align*}
\note{Le signe des trois dernières lignes est peu formé ; vu les pages précédentes on attendrait $a'-a$, etc., mais la page porte plutôt $+$.}

$\boxed{\lambda, \mu, \nu}$ correspondant aux trois directions d'arêtes

$\underline{\lambda\ \mu\ \nu}$

\drawing{Trois petits segments verticaux parallèles et un quatrième oblique, figurant les directions d'arêtes.}

\begin{align*}
a + a' &= \lambda\\
b + b' &= \mu\\
c + c' &= \nu
\end{align*}
\[
\varphi_a + \varphi_{a'} = 2(b + b' + c + c') = 2(\mu + \nu)
\]
\note{Les indices des deux $\varphi$ sont à peine formés ; la lecture $\varphi_a + \varphi_{a'}$ est suggérée par la p. 216.}

\end{document}
